\chapter{商业银行：信用中介与流动性创造}
\label{chap:Commercial_Banks_Credit_Intermediation_and_Liquidity_Creation}

\begin{openingcase}{为什么溶剂银行仍然可能面临现金短缺的问题？}

假设一家银行的资产为人民币1000亿。其大部分资产都是向家庭和企业的贷款，此外还有几年后才会到期债券的持有。在负债方面，该银行的大部分资金来自客户可以随时提取或短时间内即可取款的存款。该银行还有能够吸收损失的权益资本。

假设银行的借款人继续还款，而银行仍然能够维持偿付，其资产价值超过负债。即便如此，突然出现的提款潮仍可能导致立即的资金需求。快速出售长期资产可能需要降价出售，而用市场资金替代存款则可能成本较高或无法获得。

同样的资产负债表结构使得银行能够在提供流动性存款的同时，也能进行非流动贷款的融资。然而，当提取资金的需求突然增加时，这种结构也会带来风险。

\end{openingcase}


\begin{tcolorbox}[
  enhanced,
  colback=FinanceBlue!6,
  colframe=FinanceBlue,
  boxrule=0.6pt,
  sharp corners,
  left=8pt,
  right=8pt,
  top=7pt,
  bottom=7pt,
  before skip=1em,
  after skip=1.6em,
  title=\textbf{核心问题},
  coltitle=white,
  colbacktitle=FinanceBlue,
  fonttitle=\bfseries,
  attach boxed title to top left={xshift=6pt,yshift=-2mm},
  boxed title style={
    boxrule=0pt,
    sharp corners,
    left=5pt,
    right=5pt,
    top=2pt,
    bottom=2pt
  }
]
\begin{enumerate}[
  leftmargin=*,
  label=\textcolor{FinanceBlue}{\textbf{\arabic*.}},
  itemsep=0.55em,
  topsep=0.2em
]


\item 商业银行如何利用其资产负债表将存款人和其他提供者的资金转化为贷款和其他金融资产呢？

\item 为什么银行能够通过信息的生成以及合同的设计来创造价值？这包括对风险的筛选、监控，以及关系型贷款等方面。

\item 期限转换、流动性创造、杠杆效应以及资金结构如何在为银行创造价值的同时，也使其面临信用、利率和流动性风险？

\end{enumerate}
\end{tcolorbox}

\section{银行资产负债表与资产负债表中介}

商业银行的核心是一种资产负债表中介：通过存款和其他负债筹资，取得贷款及其他金融资产，并持有权益资本吸收损失。

\subsection{银行资产负债表：资产、融资与交易}

银行的资产负债表显示了其所持有的资产以及这些资产的融资方式。商业银行通常持有贷款、证券、储备金以及其他流动性资产。这些资产是通过存款、批发融资以及其他债务以及权益资本来融资的。

资产负债表的经济含义，体现在各类头寸的构成和契约特征中。贷款使银行承受借款人经营表现的影响。证券可提供收益、流动性和抵押品，准备金及其他流动资产支持支付结算和短期流动性需求。融资端的存款与批发借款，在期限、稳定性、重新定价和提款条件上不同。权益在债权人请求权受损之前吸收损失。

\begin{table}[H]

\centering

\caption{简化的商业银行资产负债表}

\label{tab:ch12-simplified-bank-balance-sheet}

\begin{tabularx}{0.82\textwidth}{@{}X X@{}}

\toprule

\textbf{资产} & \textbf{负债与权益} \\

\midrule

向家庭和企业的贷款 & 交易和储蓄存款 \\

政府债券及其他证券 & 定期存款和定期存单 \\

储备及其他流动性资产 & 批发和市场融资 \\

& 权益资本 \\

\bottomrule

\end{tabularx}

\end{table}

商业银行通过持有对借款人和其他发行者的请求权来进行中介业务，同时向储户和其他资金提供者发行自己的负债。双方的资产和负债在到期时间、流动性、风险和重新定价特性方面可能存在差异。双方资产的管理是共同进行的：贷款决策会影响银行资产的规模和构成，而存款流动、批发融资和权益则会影响这些资产的融资方式以及银行的可用流动性和损失吸收能力。

银行交易直接体现这些相互作用。银行发放贷款，并将款项记入借款人在本行的存款账户时，同时记录一笔新增贷款资产和相应存款负债。

\begin{example}{从贷款发放到结算}{loan-origination-balance-sheet}

在发放新贷款之前，银行拥有以下简化的资产负债表，单位以百万人民币计：

\begin{center}

\begin{tabular}{
@{\hspace{0.8em}}
p{0.25\textwidth} r
@{\hspace{2.0em}}
p{0.24\textwidth} r
@{\hspace{2.5em}}
}

\toprule

\multicolumn{2}{c}{\textbf{资产}} &
\multicolumn{2}{c}{\textbf{负债与权益}} \\

\cmidrule(lr){1-2}\cmidrule(lr){3-4}

贷款 & 70 & 存款 & 82 \\

证券 & 15 & 批发融资 & 8 \\

储备金 & 15 & 权益 & 10 \\

\midrule

\textbf{总计} & \textbf{100} &
\textbf{总计} & \textbf{100} \\

\bottomrule

\end{tabular}

\end{center}

银行提供了新的人民币约500万元的贷款，并将资金存入借款人的存款账户中。在申请过程中，

\[
\Delta L=+5,
\qquad
\Delta D=+5.
\]

总资产和总负债加上权益均增加至约1.05亿元人民币。权益在初始时保持不变。

假设借款人随后将人民币~500~万元存款转入另一家银行的账户。如果通过准备金结算，且没有抵销性流入，贷款发放银行的存款负债和准备金均减少~500~万元，贷款仍保留在资产端。

\end{example}

贷款发放与支付结算是不同的资产负债表事件。发放贷款时创造存款，本身并不提供此后资金离开银行所需的结算资产。银行遭遇存款净流出时，可能需要动用流动资产、获得替代融资，或调整其他头寸。

\subsection{银行资本、杠杆与资产负债表约束}

在这一资产负债表结构中，权益资本是满足负债后对银行资产的剩余请求权。资产价值下降先减少权益，之后才影响债权人请求权。如果损失足以耗尽资本缓冲，债权人就可能承受损失。

一种简单的资产负债表杠杆指标是资产权益比~$A/E$。杠杆越高，给定权益支持的资产规模越大，权益对资产价值变化也越敏感。负债不变时，资产损失反映为权益减少，因此资产较小比例的下降，也可能造成权益大得多的比例降幅。

\begin{example}{利用与损失吸收}{bank-leverage-loss}

该银行的资产为人民币约1000亿元，负债为人民币约920亿元，权益为人民币约80亿元。其资产与权益的比率为

\[
100\div 8=12.5.
\]

假设资产价值下降2\%，即约20亿元人民币，而负债保持不变。权益则从约80亿元人民币下降到约60亿元人民币，减少了25\%。

这个例子说明了杠杆效应如何使得权益对资产价值的变化更加敏感，而总资产则相对不那么敏感。

\end{example}

在银行资产负债表扩张的过程中，有几个限制因素需要考虑。额外的贷款或其他资产增长需要资金支持，同时也会增加资本和流动性的需求。审慎的监管要求还会进一步限制杠杆率、风险加权暴露、流动性状况以及风险集中度。究竟哪个限制因素最为重要，则取决于银行的资产负债表结构以及市场状况。

这些限制因素影响着资产负债表的规模和构成。扩大贷款规模会增加信贷风险，可能需要额外的资金和资本支持；而持有更多流动性资产虽然能增强流动性能力，但可能会降低预期回报。过度依赖短期批发融资可以提供灵活的融资方式，但同时也会增加再融资的风险。

\begin{tcolorbox}[
  enhanced,
  breakable,
  colback=white,
  colframe=FinanceBlue,
  boxrule=0.7pt,
  sharp corners,
  borderline west={2.2pt}{0pt}{FinanceBlue},
  title=\textbf{财务展望：资本与流动性发挥着不同的作用},
  coltitle=white,
  colbacktitle=FinanceBlue,
  fonttitle=\bfseries,
  before skip=1.2em,
  after skip=1.2em
]

当银行资产的价值下降时，资本能够吸收这些损失。流动性则确保了在支付和提取发生的时候能够顺利进行。

银行可能权益为正，却没有足够流动资源应付即时资金流出。反之，银行也可能持有大量流动资产，但损失已大到削弱其资本状况。

资本与流动性之间存在关联，但它们分别涉及资产负债表中的不同脆弱性方面。

\end{tcolorbox}



\section{信用评估与贷款决策}

贷款通过一种信贷流程进入商业银行的资产负债表，该流程将借款人的信息转化为贷款决策。商业银行利用与借款人相关的信息以及公开可用的信息来评估信用质量、设定信用限额和贷款条件，并在情况发生变化时重新评估风险敞口。因此，信贷分配是一个持续的过程，而不是在贷款发起时的单一决策。

在银行信贷关系的各个阶段，主要的信息需求以及相应的贷款决策都在表~\ref{tab:ch12-credit-process}中有所总结。

\begin{table}[H]
\centering
\caption{关于银行信贷关系的信息与决策}
\label{tab:ch12-credit-process}
\begin{tabularx}{0.90\textwidth}{@{}p{0.22\textwidth} X X@{}}
\toprule
\textbf{阶段} & \textbf{信息重点} & \textbf{银行决策} \\
\midrule
初步评估
& 偿还能力、财务状况、现金流稳定性、抵押品以及经营状况
& 批准、拒绝或限制拟定的发布方式。 \\
合同设计
& 资金的预期使用方式、还款模式、不确定性因素以及银行所能提供的保护措施。
& 设定金额、到期日、还款条款、定价、抵押品以及契约限制条款 \\
监控与更新
& 支付行为、账户活动情况、最新的财务信息以及业务状况的变化。
& 维护、扩展、强化、重构或终止相关内容的展示。 \\
\bottomrule
\end{tabularx}
\end{table}

\subsection{借款人信息与信用评估}

在信用评估中，会考虑借款人的履行合同义务的能力以及所承担风险的大小。这一评估基于多种信息来源进行，每一种信息来源都提供了关于信用质量不同方面的证据。

财务报表提供了关于盈利能力、杠杆率、流动性和资产负债表结构的信息。现金流记录有助于评估运营收入是否能够支持债务偿还，而还款历史和账户活动则提供了有关还款行为的证据。行业状况、客户集中度、所有权结构以及贷款的目的等因素，也会进一步影响预期的还款能力。如果合同规定的付款未能按约定履行，抵押品则可以提供额外的回收来源。

银行通过财务分析、评分系统、统计模型以及专业判断来整合这些信息。每种方法的重要性因借款人和贷款类型而异。标准化的信息有助于进行更公平的评估，而当还款能力取决于那些仅凭财务比率难以衡量的因素时，关于借款人的具体信息就变得更为重要了。

这些信息随后会被转化为内部的信用评估结果，这种评估可以采用评级、风险等级或其他银行在信贷过程中使用的分类方式来表示。不过，这种评估并不能机械地决定是否给予贷款。它只是为关于信贷额度、合同条款以及后续监控的强度和方式等决策提供参考依据而已。

\begin{example}{从借款人信息到信用评估}{borrower-information-credit-assessment}

有两家公司申请了类似的营运资金支持计划。两家公司的年度销售额均约为2000万元人民币，其杠杆比率也相当相似。

A公司能够产生稳定的运营现金流，拥有足够的债务偿还能力，并且其还款记录一直都很稳定。而B公司则面临更大的现金流波动风险，其履行债务义务的能力较为有限，而且有延迟付款的历史记录。

银行通常会认为公司A更具信用价值，因此会给其更高的内部信用评级，同时违约概率也会较低。在其他条件相同的情况下，公司A将能够获得更高的信贷额度、更低的贷款利差，以及相对较宽松的抵押品要求。

\end{example}



\subsection{贷款条款、信用限额与风险暴露}

信用评估随后会被转化为贷款决策。银行会设定信用额度和其他贷款条款，以确定可以承担的风险规模以及这种风险随时间的变化情况。主要的合同条款及其对银行信用风险的影响在表~\ref{tab:ch12-loan-terms-exposure}中进行了总结。这些条款共同决定了银行对借款人的风险敞口的规模、久期和条件。

\begin{table}[htbp]

\centering

\caption{选定的贷款条款及其对银行风险敞口的影响}

\label{tab:ch12-loan-terms-exposure}

\begin{tabularx}{0.90\textwidth}{@{}p{0.26\textwidth} X@{}}

\toprule

\textbf{贷款期限} &
\textbf{曝光影响} \\

\midrule

贷款金额或信用额度 &
规定了合同中允许使用的最大信用额度 \\

到期日 &
设定了需要偿还或续签的期限。\\

偿还结构 &
规定了存续本金随时间的变化情况 \\

利率结构 &  
决定了是否以及以何种方式重新设定合同利率 \\

抵押品和契约限制条款 &
通过抵押资产和对借款人行为的限制或要求来提供合同保护 \\

绘图与更新条件 &  
规定在何时可以提取额外资金，以及信用额度是否会在当前期限之后继续有效 \\

\bottomrule

\end{tabularx}

\end{table}

偿还和提款条款决定了在授予信贷后，存续的敞口会如何演变。用$B_t$表示第$t$期末的存续本金，$D_t$表示该期间的新提款额，$R_t$表示本金偿还额。一个简单的平衡关系为：

\begin{equation}
B_t=B_{t-1}+D_t-R_t.
\label{eq:ch12-loan-balance-dynamics}
\end{equation}

对于定期贷款，本金通常在借款时一次性支付，之后随着还款的进行而逐渐减少。而循环信贷机制则有所不同：借款人可以在合同约定的信用额度内进行提款、还款，并继续提款。因此，信用额度与当前存续的本金数额是分开的。

这一区别形成两个相关的风险暴露指标。\emph{已融资风险暴露}是已经发放且仍未偿还的金额；\emph{未提用承诺}是承诺额度中尚未使用的部分，借款人按契约进一步提款后，可能转化为已融资风险暴露。因此，即使契约信用限额不变，使用率提高也可能增加银行的已融资风险暴露。

\begin{example}{循环信贷融资与银行风险敞口}{revolving-credit-exposure}

一家公司获得了为期一年的循环信贷额度，额度为人民币约1000万元。年初时，有人民币约400万元尚未偿还。在下一个季度内，该公司又借入了人民币约300万元，并偿还了人民币约100万元。使用公式~\ref{eq:ch12-loan-balance-dynamics}来计算。

\[
B_t=4+3-1=6
\]

百万人民币。

因此，该银行拥有约600万的人民币资金敞口，而根据这项融资安排，还有约400万的人民币资金并未被使用。如果借款人日后能够使用掉这部分未使用的资金，那么资金敞口将会增加，但合同规定的信贷限额则不会发生变化。

\end{example}

贷款条款确定了初始的敞口规模，但一旦信贷关系建立后，借款人的状况可能会发生变化。因此，新的信息可能会导致银行重新评估该敞口，进一步考虑是否需要追加贷款，以及是否应该继续维持这一贷款关系。

\subsection{监督、续贷与信贷调整}

通过监控，银行能够利用还款情况、账户活动、财务信息、抵押品价值、契约限制条款的遵守情况以及业务状况的变化来更新其评估。契约限制条款指的是借款人是否继续满足贷款合同中规定的限制和要求。

这些指标的变动可能会引发信用审查。未能按时支付、经营现金流的持续疲软，或者严重的契约限制条款违反，都可能表明状况恶化；而更良好的支付记录或财务状况则可能支持现有信用的持续。进行此次审查的目的是为了确定银行当前的风险敞度是否与其对借款人的最新评估一致。

续签决策是一个重要的节点。银行在决定是否继续、扩大、减少或终止这一信贷安排时，可以将当前的信息与借款人在整个贷款期限内的表现结合起来考虑。因此，即使没有发生违约情况，通过借贷关系积累的信息仍然会影响后续的信贷决策。

信用调整取决于重新评估的结果以及合同中规定的权利条款。在贷款期限内，银行可以在允许的情况下限制借款人的额外提取行为，或者要求借款人采取特定的保护措施。在贷款到期时，银行可以重新确定信用限额以及其他合同条款。因此，通过定期监测和重新评估，银行能够将对借款人状况变化的影响转化为对银行风险敞口的变化调整。

\begin{example}{从监控到信用调整}{monitoring-credit-adjustment}

某公司拥有高达800万人民币的循环信贷额度。销售额基本保持稳定，但经营现金流入减少，客户应收账款增加，因此信用额度的使用比例从40\%上升到了85\%。 

这些发展本身并不意味着默认情况的发生，但综合来看，它们可能会引发对信贷状况的审查。银行可以重新评估现金流的压力来源，更新对借款人的评估，并重新考虑现有的风险状况。

如果银行认为这种状况只是暂时的，那么银行可以继续维持该贷款安排。但如果情况持续恶化，根据合同条款，银行可能会限制进一步的提款行为，或者要求采取额外的保护措施。在贷款期限到期时，银行可能会降低信用额度，甚至拒绝继续提供全额贷款。

\end{example}

因此，信贷决策不仅会影响银行资产的构成，还会影响到为这些资产提供资金支持所需的资金数额。这一因素直接关系到银行的融资结构和流动性管理问题。

\section{银行融资与流动性创造}

商业银行通常通过较短期限且流动性较高的负债来为长期投资提供资金。创造流动性取决于银行的融资结构以及其应对现金流出需求的能力。

\subsection{存款融资与批发融资}

银行融资的构成非常重要，因为不同的融资来源在提取权、到期时间、重新定价以及集中程度方面存在差异。商业银行主要依赖存款以及各种形式的批发融资，而这些资金来源对利率变化、市场状况以及信心的变化反应各不相同。

存款是商业银行负债，但不同产品和客户的资金行为不同。交易性存款与储蓄存款可能在较短通知期内支取，定期存款则通常提供期限更明确的融资。稳定性还取决于存款人集中度、存款保险覆盖、利率敏感性及对银行风险的判断。零售存款往往更分散，大额企业或机构余额则可能形成更集中的融资风险暴露。

批发融资包括银行间借款、定期存单、担保借款以及其他基于市场的负债。当资产负债表增长或支付流需要额外的资金时，批发融资可以补充存款资金。其有用性不仅取决于成本，还取决于期限和续期条件：短期融资需要更频繁的再融资，而当市场状况恶化时，这类融资可能会变得不那么可用。

\begin{table}[H]

\centering

\caption{资金来源及选定的流动性特征}

\label{tab:ch12-bank-funding-characteristics}

\begin{tabularx}{0.92\textwidth}{@{}p{0.23\textwidth} X X@{}}

\toprule

\textbf{资金来源} &
\textbf{合同特征} &
\textbf{对银行融资的影响} \\

\midrule

交易和储蓄存款&通常可以在短时间内提供给客户使用，具体详情请参阅产品条款说明。&稳定性取决于客户结构、集中度、利率以及市场信心等因素。\\

定期存款 &  
资金在指定的到期日或通知期限内保持不变 &  
在到期前，资金的供应较为可预测；而随着到期日的临近，就需要做出续存决策了 \\

定期存款 &  
可交易或不可交易的银行负债，具有特定的到期日 &  
融资成本和滚动投资取决于市场条件和投资者需求 \\

银行间借款 &
从其他金融机构获得的短期借款 &
提供灵活的融资方式，但可能需要频繁进行再融资操作 \\

抵押品融资市场融资 &
以抵押品为担保的借款 &
资金供应取决于可用的抵押品和融资条件 \\

\bottomrule

\end{tabularx}

\end{table}

融资结构与流动性创造紧密相关。银行可以用为持有人提供较高流动性的负债，融资相对缺乏流动性的资产。维持这种结构，需要足够的流动性能力，以应付提款、到期负债及其他到期现金流出。

\subsection{流动性能力与资产负债表调整}

许多银行的负债要求在银行资产到期之前进行结算。银行的流动性能力取决于在相关时间范围内能够转化为结算资金的资源。这些资源可以包括现金和储备金、可出售资产的收益，以及额外的借款能力。

一个简单的表示方式如下：保留LaTeX命令、公式、交叉引用以及数字的内容。

\begin{equation}
LC_t=C_t+M_t+F_t^{\mathrm{avail}},
\label{eq:ch12-liquidity-capacity}
\end{equation}

其中，$C_t$表示立即可用的现金和储备金，$M_t$表示在一段时间内可以从可销售资产中合理获得的现金，而$F_t^{\mathrm{avail}}$则表示可以合理获得的额外资金。为了避免重复计算，当抵押品或借款能力依赖于相同的资产时，应调整这些金额。因此，流动性能力取决于时间范围以及必须筹集资金的状况。

从某项资产中获得的现金金额可能与其报告的市场价值有所不同。交易成本、结算时间、价格变动以及抵押品使用的限制都会影响到资产销售所得的资金。同样，借款能力也取决于抵押品的可用性和价值、交易对手的借款意愿以及当前的融资条件。

当出现现金流出时，银行可以调整其资产负债表中的各项指标。它可以使用储备金、出售可交易的证券，或者通过抵押符合条件的资产来获得有担保的借款，同时还可以获得其他替代性融资，或者减少新资产的增长。这些应对措施通过不同的渠道来满足流动性需求，并会在之后产生不同的资产负债表状况。

\begin{example}{遇到存款流出情况}{meeting-deposit-outflow}

一家银行拥有约60亿元人民币的储备金，约80亿元人民币的可用证券尚未作为抵押品，还有约40亿元人民币的担保借款能力，这些能力由其他符合条件的资产来支持。该银行出现了约100亿元人民币的净存款外流。

一种可能的调整是使用60亿元人民币的储备金，并获得40亿元人民币的担保借款。另一种方法是使用储备金出售40亿元人民币的可交易证券。这两种方案都能提供所需的流动性，但它们对资产负债表的影响却有所不同。

第一种应对方式会增加抵押品所对应的负债，同时也会加重资产负担。第二种应对方式则会减少银行持有的可交易证券。因此，流动性管理不仅涉及满足当前的资金流出需求，还涉及到之后仍然可用的流动性资源。

\end{example}

流动性能力决定了银行能够多快地应对当前的现金流出需求。如果流出金额较大或持续时间较长，那么银行就需要进行更大的资产负债表调整。

\subsection{融资流出与流动性调整}

当存款从银行流出，或者到期需要偿还的负债在没有新资金补充的情况下被偿还时，就会出现资金流出。这种资金流出会减少银行的资金供应，从而产生现金需求。除非有资产端的现金流入或新的资金来源来抵消这种需求。

用 $O_t$ 表示在较短时间内发生的资金流出，$I_t$ 表示来自到期或出售资产的现金收入，而 $F_t$ 则表示在同一时期获得的新的资金。现金和储备金的变化可以表示为：

\begin{equation}
\Delta C_t=I_t+F_t-O_t.
\label{eq:ch12-liquidity-adjustment}
\end{equation}

该关系图展示了资金流出是如何影响银行的流动性状况的。当资产端的现金流入和新融资不足以弥补资金流出时，现金和储备金就会减少。而当替代资源充足时，银行就可以在不耗尽现有现金储备的情况下应对资金流出问题。

所需的调整不仅取决于资金流出金额，还取决于其时机和集中程度。对于大型批发负债的不再续期来说，可能需要立即进行替代，而如果将相同金额的提取分散到多个存款人身上，则情况会有所不同。相关的问题是，在可用的时间范围内，需要替换多少资金。

预期回报

因此，资金流出会改变银行的流动性水平以及资产负债表的结构。更为显著或持续性的调整还会影响到银行的融资成本、资产持有情况，以及银行风险的其他方面。

\section{银行综合风险管理}

银行风险是通过共同的资产负债表相互影响的。信贷质量、利率或融资方面的冲击可能会影响到企业的收益、资本和流动性，而由此产生的资产负债表调整则可能会带来更多的风险。

\subsection{贷款损失、收益与资本}

贷款投资组合的恶化可能会降低银行的收益，并通过留存收益削弱其资本状况。这一过程从确认信用损失开始，最终影响到当前的收入和权益。

贷款损失拨备费用是在适用会计框架下确认信用损失的一项费用。计提拨备减少当期收益，相应贷款损失准备减少贷款净账面金额。此后若用已有准备核销贷款，贷款总额和准备同时减少；损失已通过较早的拨备费用影响收益。

简化的银行收入报表可以这样编写：

\begin{equation}
\textit{NI}=\textit{NII}+\textit{NIR}-\textit{OE}-P_L-T,
\label{eq:ch12-bank-net-income}
\end{equation}

其中，$\textit{NI}$表示净收入，$\textit{NII}$表示净利息收入，$\textit{NIR}$表示非利息收入，$\textit{OE}$表示运营费用，$P_L$表示贷款损失准备费用，而$T$则表示税费。该方程揭示了准备费用对银行收益的直接影响。

如果股利和其他分配用$D$表示，而该期间发行的新权益用$\textit{NE}$表示，那么账面权益的变动大致如下：

\begin{equation}
E_1=E_0+\textit{NI}-D+\textit{NE},
\label{eq:ch12-equity-evolution}
\end{equation}
除了对权益的其他直接调整之外，这两个关系展示了传导机制：

\[
\begin{array}{c}
\text{\emph{损益表}} \\[2pt]
\text{贷款恶化}
\rightarrow
\text{相关费用条款}
\rightarrow
\text{较低的净收入}
\\[5pt]
\Downarrow
\\[-1pt]
\text{\emph{资产负债表}} \\[2pt]
\text{降低留存收益}
\rightarrow
\text{较弱的权益状况}
\end{array}
\]

\begin{example}{资源配置、收益以及资本}{provisioning-earnings-capital}

一家银行年初拥有6亿人民币的权益。在扣除准备金和税费后，其实际收益为1.2亿人民币。该银行记录的贷款损失准备金额为7000万元人民币，同时支付了2000万元人民币的股利。忽略税费和其他权益调整因素。

净收入为5000万人民币，年末的权益为

\[
600+50-20=630
\]

百万人民币。

如果在相同假设下，准备金费用上升至1.3亿元人民币，那么净利润将变为负1,000万元人民币，年末权益将降至5.7亿元人民币。因此，更大的已确认贷款损失会减少收益，而在其他条件不变的情况下，银行的权益将不足以吸收后续的损失。

\end{example}

\subsection{重新定价与资产负债错配}

利率变化通过收入渠道和估值渠道对银行产生影响。这种影响取决于资产和负债的利率何时重新设定，以及它们的经济价值对市场收益率变化的敏感程度。

让 $\textit{RSA}$ 和 $\textit{RSL}$ 分别表示在给定时期内需要重新估价的利率敏感型资产和负债的金额。对于某个特定的资产负债表状况，净利息收入的大致变化为

\begin{equation}
\Delta \textit{NII}\approx \textit{RSA}\cdot\Delta r_A
-RSL\cdot\Delta r_L,
\end{equation}
其中，$\Delta r_A$和$\Delta r_L$分别表示相应资产和负债的变动幅度。如果这些变动幅度大致相同，即$\Delta r_A\approx\Delta r_L\approx\Delta r$，那么表达式就可以简化为

\begin{equation}
\Delta \textit{NII}\approx (\textit{RSA}-\textit{RSL})\Delta r
=G\Delta r.
\label{eq:ch12-net-interest-income-gap}
\end{equation}

重新定价缺口~$G$~衡量利率敏感性资产与负债之差。正缺口意味着利率上升通常提高净利息收入，负缺口则相反。

这一近似刻意作了简化。存款利率对市场利率的调整可能不到一比一，贷款可能仅在约定日期重设利率，或参照不同基准，客户行为也可能改变未偿余额。因此，重新定价缺口只是短期收入敏感性的初步近似，并非利率风险暴露的完整度量。

经济价值分析考察错配的另一维度。即使近期现金流不变，固定利率资产与负债对市场收益率变化的敏感性也可能不同。设~$D_A$~和~$D_L$~分别表示资产和负债的有效久期，$A$~和~$L$~表示其当前经济价值。久期缺口为

\begin{equation}
D_{GAP}=D_A-\frac{L}{A}D_L.
\label{eq:ch12-duration-gap}
\end{equation}

有效久期衡量的是价值对收益率微小变化的近似百分比敏感性。对于市场收益率的微小平行变化，

\begin{equation}
\Delta E\approx -A\,D_{GAP}\,\Delta y,
\label{eq:ch12-duration-gap-equity}
\end{equation}

其中，$\Delta E$~表示由此产生的权益经济价值变化。因此，在上述近似条件下，正久期缺口意味着收益率上升可能降低权益经济价值。

\begin{tcolorbox}[
  enhanced,
  breakable,
  colback=white,
  colframe=FinanceBlue,
  boxrule=0.7pt,
  sharp corners,
  borderline west={2.2pt}{0pt}{FinanceBlue},
  title=\textbf{财务展望：收入效应与经济价值效应可能存在差异},
  coltitle=white,
  colbacktitle=FinanceBlue,
  fonttitle=\bfseries,
  before skip=1.2em,
  after skip=1.2em
]

同样的利率变动可能会以不同的方式影响银行的收入和经济价值。例如，利率上升可能会增加短期的净利息收入，但同时会降低权益的经济价值。

\end{tcolorbox}

\subsection{风险通过资产负债表的相互作用}

银行风险是相互影响的，因为信用损失、利率波动以及资金供应方面的冲击都会对同样的资产负债表产生影响。

某一领域的冲击可能改变另一领域的风险暴露。贷款损失可能削弱收益与资本；利率变化可能同时改变资产价值与融资成本；融资流出可能要求出售资产或借入替代资金。这些应对又可能影响抵押品可得性、已实现损益、融资成本，以及银行吸收进一步冲击的能力。

这些相互作用的强度取决于资产和负债的构成、融资期限、抵押品的可用性、套期保值措施，以及银行的初始资本和流动性状况。

\begin{figure}[H]
\centering

\begin{tikzpicture}[
  >=Latex,
  box/.style={
    draw,
    rounded corners,
    align=center,
    text width=8cm,
    minimum height=0.95cm,
    inner sep=5pt
  },
  main/.style={
    draw=FinanceBlue,
    fill=FinanceBlue!8,
    rounded corners,
    align=center,
    text width=9.6cm,
    minimum height=1.0cm,
    inner sep=5pt
  },
  flow/.style={
    -Latex,
    line width=0.8pt,
    FinanceBlue
  }
]

\node[box] (shock) at (0,3.0)
{信用、利率或融资冲击};

\node[main] (bs) at (0,1.3)
{资产负债表对现金流和价值的影響};

\node[box] (capacity) at (0,-0.4)
{收益、流动性与资本};

\node[box] (adjustment) at (0,-2.1)
{资产负债表反应};

\draw[flow] (shock) -- (bs);
\draw[flow] (bs) -- (capacity);
\draw[flow] (capacity) -- (adjustment);

\draw[flow]
(adjustment.east)
.. controls +(2.0,0) and +(2.0,0) ..
(bs.east);

\end{tikzpicture}

\caption{商业银行内的风险传递与资产负债表调整}
\label{fig:ch12-integrated-risk-interactions}
\end{figure}
资产负债表调整由此可能改变银行随后面临的信用、利率和融资风险暴露。

\subsection[资本应对与资产负债表调整]
{资本应对与资产负债表调整}

信用损失会降低账面权益，而资产快速增长如果权益增长不成正比，则可能会降低权益与资产的比例。银行可以通过保留收益或发行新股权来加强资本状况，从而减少资产基础，或者结合这两种措施来应对这一问题。一种简单、未加权的资本衡量方法就是……

\begin{equation}
k=\frac{E}{A},
\end{equation}
其中，$E$表示账面权益，而$A$则表示总资产。  
在持有资产不变的情况下，权益的增加会提高这一比率；而在保持权益不变的情况下，资产的减少会产生同样的效果。

通过保留收益、减少分配或发行新股份，可以增强权益。资产调整可以包括减缓新的贷款发放，让资产在无需替换的情况下到期，或者出售贷款和证券。这些选择在时机和经济成本上有所不同，同时也会改变未来的收益和资产负债表的构成。

假设一家银行的初始权益为$E_0$，资产为$A_0$，在确认资产损失$L_s$时，负债保持不变，那么……

\begin{equation}
E_1=E_0-L_s,
\label{eq:ch12-capital-after-loss}
\end{equation}
而破产后的资产则……

\begin{equation}
A_1=A_0-L_s.
\end{equation}

如果没有新的权益被筹集，而银行寻求一个目标比率$k^{*}$，那么与这一目标相一致的资产水平是

\begin{equation}
A^{*}=\frac{E_1}{k^{*}}.
\label{eq:ch12-assets-consistent-capital-ratio}
\end{equation}

当 $A_1 > A^{*}$ 时，差值 $A_1 - A^{*}$ 表示在这些简化假设条件下所需的资产减少量。

\begin{example}{恢复资本比率的两条途径}{capital-response-paths}

一家银行的资产为人民币约1000亿元，权益为人民币约80亿元。该公司确认了约20亿元人民币的资产损失，因此目前资产为人民币约980亿元，权益为人民币约60亿元。该公司的权益与资产比率从8\%下降到了约6.12\%。

假设管理层希望恢复8\%的比例。

一种方法是发行新的权益。如果募集资金仍然留在资产负债表上，那么发行约20亿元人民币的权益将使得资产总额达到约1000亿元人民币，而权益则达到约80亿元人民币：

\[
\frac{8}{100}=8\%.
\]

第二种方法是不增加新的权益，而是减少资产。在权益仍为约60亿元人民币的情况下，资产水平与8\%的比例相符合。

\[
\frac{6}{0.08}=75
\]

十亿人民币。

因此，该银行需要将资产减少约230亿元人民币，以弥补亏损后的资产规模。这两种解决方案虽然保持了相同的比率，但资产负债表中的各项指标却有所不同。

\end{example}

因此，资本结构调整不仅会影响资本比率，还会影响到银行的未来收益、流动性、融资需求以及风险暴露情况。

%========================================================
% Summary
%========================================================

\TailModule{summary}{本章小结}

商业银行是一种资产负债表中介机构。它持有贷款、证券、储备金以及其他资产，同时通过存款、批发负债和权益来为这些资产提供资金。贷款的发放可以同时产生贷款资产和存款负债；而后续的支付则可以将存款和储备金在银行之间转移。在债权人的请求权受损之前，权益会先吸收资产损失；而且，杠杆效应会增加权益对资产价值变化的敏感性。

银行信贷决策是随着时间而变化的。借款人的相关信息有助于评估信用状况，而贷款条款和额度则决定了银行的风险规模和结构。在贷款发放后，持续的监控能够确保获得最新的信息；而当借款人情况发生变化时，贷款的续期或合同调整也会对银行的风险产生影响。

银行融资将存款与具有不同到期日、集中程度、重新定价条件和提取条件的批发负债相结合。当相对流动性较高的负债用于资助流动性较低或期限较长的资产时，就会产生流动性创造。因此，流动性管理取决于可用于满足现金流出需求的资源，以及当融资条件发生变化时所使用的资产负债表调整措施。

银行风险通过共同的资产负债表相互关联。贷款恶化可能会减少收益和资本，利率波动可能会影响净利息收入和经济价值，而资金流出则可能会消耗流动性资源或需要新的融资方式。对一种压力的应对可能会改变对另一种压力的风险承受能力。权益可以通过留存收益或新的权益来重建，而资产调整则可能会改变资产负债表的规模和构成，从而影响未来的收益、流动性和风险。

%========================================================
% Mind Map
%========================================================

\TailModule{mindmap}{思维导图}

\begin{figure}[H]
\centering
\begin{tikzpicture}[
  >=Latex,
  center/.style={
    draw=FinanceBlue,
    fill=FinanceBlue!8,
    rounded corners,
    align=center,
    text width=3.2cm,
    minimum height=1.0cm
  },
  box/.style={
    draw,
    rounded corners,
    align=center,
    text width=3.3cm,
    minimum height=0.95cm
  },
  flow/.style={-Latex, line width=0.8pt, FinanceBlue}
]

\node[center] (bank) at (0,0) {商业银行\\资产负债表};

\node[box] (balance) at (0,3.0)
{资产负债表中介、\\资本，以及杠杆率};

\node[box] (credit) at (-5.0,0.9)
{引用评估，\\贷款条款，以及监控措施。};

\node[box] (funding) at (5.0,0.9)
{存款、资金供应以及\\流动性管理};

\node[box] (risk) at (0,-3.0)
{风险相互作用\\资产负债表调整};

\draw[flow] (balance) -- (bank);
\draw[flow] (credit) -- (bank);
\draw[flow] (funding) -- (bank);
\draw[flow] (bank) -- (risk);

\end{tikzpicture}
\caption*{商业银行通过统一的资产负债表来管理信贷决策、资金供应、流动性和资本问题。}
\end{figure}

%========================================================
% Key Terms
%========================================================

\TailModule{keyterms}{关键术语}

商业银行（Commercial bank）；资产负债表中介（balance-sheet intermediation）；存款（deposits）；批发融资（wholesale funding）；权益资本（equity capital）；杠杆（leverage）；信用评估（credit assessment）；信用限额（credit limit）；已融资风险暴露（funded exposure）；未提用承诺（undrawn commitment）；贷款监督（loan monitoring）；贷款续期（loan renewal）；融资期限（funding maturity）；流动性创造（liquidity creation）；流动性能力（liquidity capacity）；融资流出（funding outflows）；净利息收入（net interest income）；重新定价缺口（repricing gap）；资产负债错配（asset--liability mismatch）；久期缺口（duration gap）；权益经济价值（economic value of equity）；资本比率（capital ratio）；资产负债表调整（balance-sheet adjustment）。

%========================================================
% Financial AI Lab
%========================================================

\TailModule{aiprompt}{金融~AI~实验室}

\textbf{实验室：对综合银行资产负债表进行分析与审计}

请让人工智能系统分析一家包含贷款、证券、储备金、存款、批发融资和权益的简化版商业银行。提供有关借款人信用审查、存款外流以及市场利率变化的信息。

将分析过程分为四个阶段进行实施。

\begin{enumerate}

\item \textbf{检查资产负债表恒等式。} 在贷款发放、支付结算、资金流出以及确认的损失发生后，确认资产等于负债加上权益。对于每一变化，确定相应的资产负债表条目。

\item \textbf{请评估信用决策的结果。}  
引入较为薄弱的借款人现金流信息，并判断是否应保持现有的信用额度。在评估还款能力、合同条款以及风险状况时，不应仅仅接受没有依据的违约概率。

\item \textbf{将流动性与资本分离开来。}  
实施存款流出政策，并确定在指定时间内能够提供现金的资源。需要注意的是，分析应区分银行的流动性能力、银行资产的经济价值以及银行的资本状况。

\item \textbf{追踪风险互动。}  
实施利率调整以及贷款损失准备金政策。分析这些调整对净利息收入、经济价值、收益和资本的影响，并确认任何拟议的资产负债表调整在内部是协调一致的。

\end{enumerate}

%========================================================
% Discussion Questions
%========================================================

\TailModule{questions}{讨论题}

\begin{enumerate}

\item 为什么商业银行要充当资产负债表中介的角色，而不是直接将储户现有的存款转移给借款人呢？

\item 如何根据借款人的具体情况来决定银行愿意提供的信贷额度以及相关的合同条款呢？

\item 为什么一个已承诺的信用额度会与银行当前的可用资金额度有所不同呢？

\item 资金集中度、债务期限以及利率敏感性如何影响着银行的流动性管理呢？

\item 如何在同一利率上升的情况下，增加短期的净利息收入，同时降低权益的经济价值？

\item 为什么为了增加流动性而进行的资产出售，仍然会通过已实现的利润或损失来影响企业的收益或资本状况呢？

\end{enumerate}

%========================================================
% Applied Problems
%========================================================

\TailModule{appprobs}{应用题}

\begin{enumerate}

\item \textbf{贷款发起与结算。}

一家银行最初拥有约8000万的人民币贷款、1500万的证券以及500万的储备金。其资金来源包括8200万的人民币存款、800万的批发融资以及1000万的权益。

\begin{enumerate}

\item 银行提供了新的400万人民币贷款，并将资金存入借款人的存款账户。请展示这一操作对资产负债表的直接影响。

\item 借款人将400万人民币的存款转存至另一家银行。如果结算导致储备减少400万人民币，请展示相应的资产负债表。

\item 如果银行希望将储备金恢复到300万人民币的水平，请提出两种能够产生所需现金的资产负债表调整方案。

\end{enumerate}

\item \textbf{信用额度与风险敞口。}

一家公司的流动资金为1200万人民币。初始资金为500万人民币。在本季度内，该公司提取了400万人民币，同时偿还了200万人民币的债务。

\begin{enumerate}
\item 计算最终的融资规模。
\item 计算在信用额度下未使用的金额。
\item 假设运营现金流减弱，而应收账款上升。请解释银行可以在不需要立即偿还全部存续余额的情况下，采取两种方式来调整风险敞口。
\item 解释为什么提高贷款利率并不是应对信用评估结果不佳的唯一办法。
\end{enumerate}

\item \textbf{资金流出与流动性能力。}

一家银行拥有70亿人民币的储备金，以及90亿人民币的证券，这些证券可以在一天内转换为85亿人民币的现金。此外，该银行还有30亿人民币的可动用抵押借款能力，并得到了其他抵押品的支持。然而，该银行出现了140亿人民币的净存款外流。

\begin{enumerate}
\item 根据方程~(\ref{eq:ch12-liquidity-capacity})计算流动性能力。
\item 所宣称的产能是否足以满足输出需求呢？
\item 如果银行使用所有储备金以及通过证券销售获得的约70亿元人民币现金，那么在上述假设下，还剩下多少流动资金呢？
\item 如果证券投资组合中的部分证券已经被用作抵押品，那么答案将会如何改变？
\end{enumerate}


\item \textbf{重新定价与经济价值风险分析。}


在接下来的一年里，这家银行拥有约850亿元人民币的利率敏感资产，以及约700亿元人民币的利率敏感负债。

\begin{enumerate}
\item 计算重新定价差距。
\item 如果使用方程~(\ref{eq:ch12-net-interest-income-gap})来计算，那么当相关利率上升80个基点时，年度净利息收入的变动幅度是多少呢？
\item 同一家银行的资产为人民币1000亿，负债为人民币920亿，资产久期为4.2，负债久期为1.8。计算其久期差距。
\item 估算权益经济价值的一阶变化，以应对收益增加50个基点的情况。
\item 解释为什么收入和经济价值指标在同一时期并不需要保持相同的增长方向。
\end{enumerate}

\item \textbf{资本响应与资产负债表调整。}

一家银行的资产为人民币约1200亿元，权益为人民币约96亿元。该银行确认了人民币约24亿元的亏损。

\begin{enumerate}
\item 计算亏损发生前后的权益与资产比率，假设亏损导致资产和权益减少的幅度相同。
\item 如果银行希望实现8\%的权益与资产比例，并且不增加任何新的权益，那么请计算出符合该比例的资产规模。
\item 计算亏损后的资产负债表所显示的资产减少情况。
\item 提出两种缓解资产负债表收缩的方法，并分别说明每种方法的局限性。
\end{enumerate}

\end{enumerate}

%========================================================
% Data Application
%========================================================

\TailModule{dataapp}{数据应用}

\textbf{中国的商业银行资产负债表、资金状况及资本结构}

参考以下翻译：

使用权威的中国来源，如中国人民银行、国家金融监督管理局，或上市商业银行发布的财务报表。选择一家银行或一个一致的银行业系列数据，收集至少五年的主要资产、存款、其他融资、单独报告的流动资产以及权益方面的数据。

构建能够回答三个问题的比率：银行的资产是如何分配的？资产负债表是如何融资的？有多少权益能够支持资产基础？有用的衡量指标可能包括贷款与总资产的比率、存款与总负债的比率、流动资产与总资产的比率，以及权益与总资产的比率。在整个样本期间使用一致的会计定义，并区分会计权益与监管资本指标。

找出一个时期，在这段时间内，资产增长、存款增长、资金构成或权益发生了显著变化。通过资产负债表的关系来分析这种变化，而不是孤立地解释某个比率。如果贷款的增长速度超过了存款的增长速度，那么需要确定是哪些负债或权益为这一差异提供了资金支持。如果权益的增长速度低于资产的增长速度，那么就需要计算由此导致的简单杠杆变化。

如果两家银行拥有可比的数据，那么请比较它们的资产组合、存款结构、对非存款资金的依赖程度以及流动资产状况。探讨这些差异如何影响其资产负债表在资金流出或利率变化情况下的反应。不要仅依据报告的流动资产来衡量一家银行的流动性能力。

请列出该机构的信息、样本周期、会计定义、相关文件，以及分析中使用的所有处理步骤。

%========================================================
% Case Analysis
%========================================================

\TailModule{caseanalysis}{案例分析}

\textbf{案例：硅谷银行——个体银行资产负债表分析}

2023年3月10日，加利福尼亚州金融保护与创新部门关闭了Silicon Valley Bank，并由联邦存款保险公司担任接管人。联邦储备系统随后的审查发现，该银行存在资产负债表快速增长、业务模式集中、严重依赖未投保的存款，以及流动性和利率风险管理方面的不足。

随着2022年市场利率的上升，SVB持有的长期固定利率证券的市场价值有所下降。与此同时，风险投资活动的减弱导致存款流入减少，而一些由科技或风险投资支持的客户则提取了存款以用于融资。2023年3月8日，SVB金融集团宣布进行资产负债表重组，其中包括出售一批可供出售的证券投资组合、确认税后亏损，以及计划进行资本筹集。

从银行资产负债表的角度来分析这一情况。需要解决四个问题：证券投资组合如何应对较高的市场利率？存款结构如何影响资金的稳定性？哪些资产或资金来源能够在所需的时间范围内提供流动性？实际发生的损失对银行的资本状况产生了什么影响？而计划中的资本筹集措施又该如何应对？

接下来，我们来探讨一下这些维度之间的相互作用。例如，为了增加流动性而进行的资产出售，可能会将市场价值的未实现下跌转化为实际损失，从而降低收益和资本水平。因此，这一事件表明，利率风险、资金流出、流动性调整以及资本等因素是如何通过相同的资产负债表相互影响的。

这里的分析仍然停留在单个银行的层面。关于系统范围内的传染效应、紧急流动性支持以及审慎监管方面的内容，将在后续章节中进行探讨。

\emph{主要来源：联邦储备系统理事会，  
《对硅谷银行的监管与监管评估》，2023年4月；存款保险局，  
2023年硅谷银行倒闭信息，  
SVB金融集团，SEC文件，2023年3月。}

%========================================================
% Further Reading
%========================================================

\TailModule{furtherread}{延伸阅读}

\begin{itemize}

\item Berger, Allen N., and Gregory F. Udell (1995),
``Relationship Lending and Lines of Credit in Small Firm Finance（小企业融资中的关系型贷款与信贷额度）,''
\emph{The Journal of Business}, 68(3), 351--381.

\item Kashyap, Anil K., Raghuram Rajan, and Jeremy C. Stein (2002),
``Banks as Liquidity Providers: An Explanation for the Coexistence of Lending and Deposit-Taking（作为流动性提供者的银行：对放贷与吸收存款并存的解释）,''
\emph{The Journal of Finance}, 57(1), 33--73.

\item Drechsler, Itamar, Alexi Savov, and Philipp Schnabl (2021),
``Banking on Deposits: Maturity Transformation without Interest Rate Risk（依托存款的银行经营：无利率风险的期限转换）,''
\emph{The Journal of Finance}, 76(3), 1091--1143.

\end{itemize}

\ChapterTail{summary,mindmap,keyterms,aiprompt,questions,appprobs,dataapp,caseanalysis,furtherread}
